
\section{Objectifs du projet}

\subsection{Objectif général}

Dans le cadre de notre formation universitaire en informatique, nous avons pour objectif général de concevoir et développer une simulation interactive en Java. Cette simulation représente un environnement de surveillance dans lequel des agents, appelés « gardiens », patrouillent une grille à deux dimensions. Ces gardiens sont capables de détecter la présence d’intrus et d’adapter leur comportement pour tenter de les intercepter. Ce projet a pour but de mettre en application les compétences acquises dans plusieurs unités d’enseignement telles que la programmation orientée objet, la modélisation, les algorithmes,  ainsi que la conception d’interfaces graphiques avec Swing.
\subsection{Objectifs spécifiques}


\subsubsection{ Modéliser un environnement en grille 2D}

Nous souhaitons tout d’abord concevoir une représentation logique de l’environnement sous la forme d’une grille bidimensionnelle. Chaque case de cette grille pourra représenter un élément : un obstacle, un agent (gardien ou intrus), ou une zone libre. Ce choix de modélisation favorise la compréhension des structures de données et permet de gérer efficacement les déplacements et interactions entre agents.

\subsubsection{Implémenter des gardiens dotés d’un comportement autonome}

Les agents gardiens auront un comportement autonome, basé sur une logique de patrouille. Ils devront pouvoir se déplacer selon un itinéraire défini, mais aussi modifier leur comportement en fonction des événements détectés (par exemple, en cas de détection d’un intrus). Cet objectif nous permet d’appliquer les principes fondamentaux de la programmation orientée objet et de la programmation événementielle.

\subsubsection{ Définir un champ de vision limité pour chaque agent}

Afin de rendre la simulation réaliste, nous définirons un champ de vision propre à chaque gardien. Ce champ sera limité par un rayon, et influencé par la présence d’obstacles. Ainsi, la détection d’un intrus ne pourra se faire que si ce dernier est visible depuis la position du gardien. Cet objectif nous introduit à la modélisation de la perception dans un espace discret.

\subsubsection{Mettre en place des algorithmes de détection et de poursuite}

Lorsque les gardiens détectent un intrus, ils doivent être capables de réagir. Pour cela, nous allons concevoir des algorithmes simples de poursuite permettant aux gardiens de rejoindre l’intrus en tenant compte de la topologie de la grille et des obstacles. Ce travail représente notre première approche des techniques d’intelligence artificielle appliquées à la navigation.

\subsubsection{ Visualiser graphiquement l’environnement et les déplacements}

Une partie essentielle de notre travail consiste à rendre visible l’état de l’environnement. Nous allons développer une interface graphique à l’aide de la bibliothèque Java Swing, permettant de visualiser la grille, les agents, les obstacles et les déplacements en temps réel. Cela nous permettra de renforcer notre maîtrise de la conception d’interfaces et d’améliorer la compréhension globale de la simulation.

\subsubsection{ Structurer le projet selon une architecture logicielle claire}

Pour garantir la qualité et la maintenabilité de notre projet, nous adopterons une architecture logicielle modulaire. Nous distinguerons notamment le modèle (représentation des données), le moteur de jeu (logique), et l’interface graphique (affichage). Cette séparation respecte les bonnes pratiques de développement logiciel et facilite l’évolution du système.

\subsubsection{ Favoriser l’extensibilité future du système}

Enfin, nous concevrons notre système de manière à pouvoir facilement ajouter de nouvelles fonctionnalités par la suite. Parmi les extensions possibles, nous envisageons :
\begin{itemize}
    \item l’ajout de nouveaux types d’agents (joueurs humains, caméras, etc.) ;
    \item des règles de jeu supplémentaires (modes multijoueur, chronomètres, objectifs) ;
    \item ou encore la génération aléatoire de cartes et d’événements.


\end{itemize}
\subsubsection{ Intégrer une gestion temporelle des événements}

Nous souhaitons enrichir notre simulation par une gestion du temps. Cela permettra d’ajouter des délais entre les actions, de simuler des comportements en fonction d’horloges internes, évènements déclenchés par minute, etc.). Ce mécanisme améliore le réalisme de la simulation et nous initie à la gestion des systèmes réactifs dans un contexte temps réel.

\subsubsection{Simuler des comportements d'intrus dynamiques}

Afin de rendre le scénario plus interactif, nous implémenterons des intrus capables de se déplacer librement et de prendre des décisions. Ces intrus pourront par exemple tenter d’éviter les gardiens, chercher des zones sûres ou atteindre un objectif spécifique. Cela permet d’introduire une forme d’intelligence adaptative rudimentaire, tout en simulant des comportements plus réalistes et complexes.

\subsubsection{ Ajouter un mode de configuration personnalisée}

Nous ajouterons la possibilité de configurer l’environnement au démarrage : taille de la grille, nombre de gardiens et d’intrus, emplacements des obstacles, rayon de vision, etc. Cette fonctionnalité favorisera la réutilisabilité du système pour différentes expériences de simulation, et nous permettra de valider notre architecture modulaire à travers divers scénarios.


Afin de rendre le scénario plus interactif, nous implémenterons des intrus capables de se déplacer librement et de prendre des décisions. Ces intrus pourront par exemple tenter d’éviter les gardiens, chercher des zones sûres ou atteindre un objectif spécifique, tout en simulant des comportements plus réalistes et complexes.
\subsubsection{Analyser les données de simulation à l’aide de statistiques graphiques}

Nous souhaitons intégrer un module d’analyse statistique des comportements simulés. L’objectif est de représenter visuellement, sous forme d’histogrammes et de courbes, des données telles que : le nombre de détections, les durées de poursuite, les trajets parcourus ou encore les fréquences d’apparition des intrus dans certaines zones. Ces visualisations nous permettront d’évaluer objectivement l’efficacité des gardiens et de comparer différentes configurations. Ce travail introduit des notions d’analyse de données et de visualisation graphique, en lien direct avec les enseignements en statistiques et en outils de visualisation.

Cet objectif nous pousse à anticiper les besoins futurs, à structurer notre code proprement, et à adopter une démarche genie logicielle rigoureuse.

\vspace{1em}

\subsubsection{Conclusion sur les objectifs spécifiques}

\vspace{0.5em}

Ainsi, chacun de ces objectifs constitue une étape motivante et essentielle qui nous engage pleinement dans un apprentissage actif, structuré et orienté vers la réussite.

\vspace{1em}

\noindent \textbf{En résumé}, l’ensemble de ces objectifs spécifiques nous permet de combiner théorie et pratique à travers un projet complet, stimulant et riche en apprentissages, parfaitement aligné avec les compétences attendues dans le domaine du génie logiciel.
